\begin{pastexam}[label={px:nsysu-111-osds-03}]{中山資工所甲組 111 學年度 第 3 題}
（問答）

(10\%) The following numbers are input into an empty binary search tree sequentially: 3, 6, 4, 5,
2, 7, 1. And, then 3 is deleted from the tree. When a nonleaf node is deleted, its inorder successor
will replace its position. Please draw the tree after the deletion of 3.
\end{pastexam}

\begin{pastsoln}
答案：刪除後根是 $4$；$4$ 的左子樹是 $2$（$2$ 的左子節點 $1$），右子樹是 $6$（$6$ 的左子節點 $5$、
右子節點 $7$）。 圖見下面右邊。

\textbf{插入}（資結書虛擬碼~9.6）：每個數從根往下比，小的往左、大的往右，走到空的位置接上。
\begin{itemize}
\item $3$ 成為根。 $6 > 3$ 接在 $3$ 的右邊。
\item $4$：$> 3$ 往右，$< 6$ 接在 $6$ 的左邊。 $5$：$> 3$、$< 6$、$> 4$，接在 $4$ 的右邊。
\item $2 < 3$ 接在 $3$ 的左邊。 $7$：$> 3$、$> 6$，接在 $6$ 的右邊。 $1$：$< 3$、$< 2$，接在
$2$ 的左邊。
\end{itemize}

\textbf{刪除 $3$}（資結書第~9.2~節）：$3$ 有兩個子節點，照題目的規定用中序後繼取代。 中序後繼是
右子樹裡最小的鍵：從 $6$ 一路往左走到 $4$，$4$ 沒有左子節點，所以後繼是 $4$。 做法：
\begin{enumerate}
\item 把 $4$ 從原來的位置拿掉。 $4$ 沒有左子節點、有右子節點 $5$，就讓 $5$ 接到 $4$ 原來的位置，
成為 $6$ 的左子節點。
\item 用 $4$ 取代 $3$ 的位置，成為新的根，左右子樹沿用 $3$ 原來的。
\end{enumerate}

\begin{center}
\begin{tikzpicture}[x=0.9cm,y=1.0cm,v/.style={circle,draw,inner sep=0pt,minimum size=6mm}]
\node at (2,0.8) {刪除前};
\node[v] (a3) at (2,0) {3};
\node[v] (a2) at (0.8,-1) {2}; \node[v] (a6) at (3.2,-1) {6};
\node[v] (a1) at (0.2,-2) {1}; \node[v] (a4) at (2.6,-2) {4}; \node[v] (a7) at (3.8,-2) {7};
\node[v] (a5) at (3.2,-3) {5};
\draw (a3)--(a2) (a3)--(a6) (a2)--(a1) (a6)--(a4) (a6)--(a7) (a4)--(a5);
\begin{scope}[xshift=6.5cm]
\node at (2,0.8) {刪除 $3$ 後};
\node[v] (b4) at (2,0) {4};
\node[v] (b2) at (0.8,-1) {2}; \node[v] (b6) at (3.2,-1) {6};
\node[v] (b1) at (0.2,-2) {1}; \node[v] (b5) at (2.6,-2) {5}; \node[v] (b7) at (3.8,-2) {7};
\draw (b4)--(b2) (b4)--(b6) (b2)--(b1) (b6)--(b5) (b6)--(b7);
\end{scope}
\end{tikzpicture}
\end{center}

核對：刪除後的中序是 $1, 2, 4, 5, 6, 7$，由小到大，仍是二元搜尋樹，而且正好是原來的鍵去掉 $3$。
常見的錯是把 $4$ 搬上去以後忘了 $5$：$5$ 要接到 $6$ 的左邊，不能跟著 $4$ 一起移走，也不能直接丟掉。
\end{pastsoln}
